\documentclass{article}
\usepackage{amsmath,amssymb}
\begin{document}
\section*{MAX-COMP-DUAL1: Exact finite OBDD acquisition dual}
\textbf{Status:} P versus NP OPEN. PROVED/EXACT for the declared finite weighted graph; computational validation for 44 test functions. No general SAT lower bound.
\paragraph{Model.} Fix a Boolean function f and variable orders as vertices. A directed edge exchanges adjacent variables, costing one plus the number of nonterminal nodes in the reduced OBDD rebuilt at the destination order. Let d(v) be the shortest-path distance from the identity order. Then for every edge u to v of cost c, d(v)-d(u) <= c. Summation gives cost(path) >= d(target)-d(source). Equality holds for a shortest path. The minimum cost to reach any minimum-node order is the minimum of d(v) over orders with globally minimal node count. This is a shortest-path dual certificate, not a proof that the cost model captures all Boolean compilers.
\paragraph{Controls.} Tested AND, PARITY, TRUE, UNSAT, connected contradictory CNF, independent-product CNF, and seeded random 3-CNFs. Enumerated all orders for n=3..6; 44 functions, 9,570 order evaluations. Independent Bellman-Ford computation matched Dijkstra distances exactly. All dual edge inequalities held (zero violations). Runtime 0.6835 seconds. This confirms the finite optimization harness, not a superpolynomial asymptotic bound.
\paragraph{Acquisition guard.} The experiment charges each adjacent-order rebuild but does not charge the preliminary truth-table computation or represent arbitrary compilation. Direct source decision, semantic oracle access, unrestricted circuit representation, and different order-finding algorithms bypass the graph. In particular the constant TRUE and UNSAT functions have zero-node OBDDs despite potentially difficult CNF encodings.
\paragraph{Prior work.} The project MAX-COMPRESSION-FRONTIER1 enumerated minimum OBDD sizes; this extends it with a formally exact costed path and dual potential. Darwiche and Marquis, A Knowledge Compilation Map, JAIR 2002, distinguish succinctness, query support and transformation costs. Wegener, Branching Programs and Binary Decision Diagrams, SIAM 2000, treats order-dependent OBDD complexity. Neither yields an unrestricted SAT decision lower bound.
\paragraph{Next gap.} Replace the artificial adjacent-order rebuild graph by a source-checkable transition graph incorporating unit propagation, resolution, RUP, semantic sharing, and implicit access. Find a dual potential valid for every allowed transition, and test oracle/direct-decision bypass. Even a valid lower bound remains restricted until arbitrary-algorithm transfer is proved.
\end{document}
